\documentclass[conference]{IEEEtran}

\usepackage{multirow}
\usepackage{float}
\usepackage{cite}
\usepackage{caption}
\usepackage{subcaption}
\usepackage{amsmath,amssymb,amsfonts,amstext}
% \usepackage{algorithmic}
\usepackage[dvipdfmx]{graphicx}
\usepackage{textcomp}
\usepackage{xcolor}
\usepackage{colortbl}
\usepackage[T1]{fontenc}
\usepackage{mdframed}
\usepackage{algorithm, algpseudocode} % texlive-science
%\usepackage[braket, qm]{qcircuit}
\usepackage{newtxtext}
\usepackage[normalem]{ulem}
\usepackage{nameref}
\usepackage{tikz}
\usetikzlibrary{arrows.meta, shapes, positioning, calc, fit}
\usepackage{quantikz}

%\newcommand*\rtofstarg{\tikz[baseline=(char.base)]{
            \node[shape=circle,double,draw,inner sep=1pt,minimum width=16pt] (char) {S};}}
\newcommand*\rtofsdgtarg{\tikz[baseline=(char.base)]{
            \node[shape=circle,double,draw,inner sep=1pt] (char) {S$^{\dagger}$};}}
\newcommand*\rtoftarg{\tikz[baseline=(char.base)]{
            \node[shape=circle,draw,inner sep=-0.5pt,fill=none] (char) {$\oplus$};}}
\newcommand*\rtofsctrl{\tikz[baseline=(char.base)]{
            \node[shape=circle,double,draw,inner sep=1pt,fill=black, minimum width=3pt] (char) {};}}
\newcommand*\rtofsctrlo{\tikz[baseline=(char.base)]{
            \node[shape=circle,double,draw,inner sep=1pt,fill=white, minimum width=3pt] (char) {};}}
%\newcommand*\rtofsctrlo{\tikz[baseline=(char.base)]{
%            \node[shape=circle,double,draw,inner sep=1pt,fill=white, minimum width=3pt] (char) {};}}
\newcommand*\dashedwire{\tikz{
            \draw[dashed] (0,0) -- (0.5,0); }}

\newcommand*\ctrlrtofs[1]{\push{\rtofsctrl}\qwx[#1]\qw}
\newcommand*\ctrlortofs[1]{\push{\rtofsctrlo}\qwx[#1]\qw}
\newcommand*\targrtofs{\push{\rtofstarg}\qw}
\newcommand*\targrtof{\push{\rtoftarg}\qw}
% \newcommand*\qcbox[1,2,3,4]{\gategroup[4,steps=2,style={dashed,rounded corners,fill=blue!20, inner xsep=2pt},background,label style={label position=below,anchor=north,yshift=-0.2cm}]{Compute}}


\makeatletter
\newcommand{\sublabel}[1]{\protected@edef\@currentlabel{\thefigure(\thesubfigure)}\label{#1}}
\makeatother
\captionsetup{subrefformat=parens}
\algrenewcommand\algorithmicindent{0.5em}
\newtheorem{definition}{Definition}
\newtheorem{example}{Example}
\newcommand{\thickhline}{%
    \noalign {\ifnum 0=`}\fi \hrule height 1pt
    \futurelet \reserved@a \@xhline
}
%\def\BibTeX{{\rm B\kern-.05em{\sc i\kern-.025em b}\kern-.08em
%    T\kern-.1667em\lower.7ex\hbox{E}\kern-.125emX}}
\newcommand{\figcaption}[1]{\def\@captype{figure}\caption{#1}}
\newcommand{\tblcaption}[1]{\def\@captype{table}\caption{#1}}

\newcommand{\Del}[1]{\textcolor{red}{\sout{#1}}}
\newcommand{\Add}[1]{\textcolor{blue}{\emph{#1}}}

\newcount\NGCcolor
\NGCcolor=1
\newcount\NGCcolortwo
\NGCcolortwo=1
%\NGCcolor=0
\ifnum\NGCcolortwo<1
  \newcommand{\blueoutt}[1]{\textcolor{blue}{#1}}  
\else
  \newcommand{\blueoutt}[1]{#1}
\fi
  
\ifnum\NGCcolor<1
  \newcommand{\redout}[1]{\textcolor{red}{#1}}
  \newcommand{\blueout}[1]{\textcolor{blue}{#1}}
  \newcommand{\greenout}[1]{\textcolor{green}{#1}}
  \newcommand{\cyannout}[1]{\textcolor{cyan}{#1}}
  \newcommand{\magentaout}[1]{\textcolor{magenta}{#1}}
  \newcommand{\yellowout}[1]{\textcolor{yellow}{\bf #1}}
  \newcommand{\algblueout}[1]{\leavevmode\color{blue}{#1}}
\else
  \newcommand{\redout}[1]{#1}
  \newcommand{\blueout}[1]{#1}
  \newcommand{\greenout}[1]{#1}
  \newcommand{\cyannout}[1]{#1}
  \newcommand{\magentaout}[1]{#1}
  \newcommand{\yellowout}[1]{#1}
  \newcommand{\algblueout}[1]{\leavevmode {#1}}
\fi
% Preamble (requires tikz, which quantikz already uses)
\newcommand{\qwiregap}[1]{%
  \gate[style={draw=none}]{#1}%
}
\begin{document}

\title{Synthesis of Phase Oracles Using CNOT+T Circuits and ESOP Minimization}
\author{
    \IEEEauthorblockN{\phantom{David Clarino}}
    \IEEEauthorblockA{\phantom{Ritsumeikan University} \\
      \phantom{dclarino@fc.ritsumei.ac.jp}}
    \and 
    \IEEEauthorblockN{\phantom{Chitranshu Arya}}
    \IEEEauthorblockA{\phantom{Netaji Subhas} \\
      \phantom{University of Technology} \\
      \phantom{chitranshu.arya.ug22@nsut.ac.in}}
    \and
    \IEEEauthorblockN{\phantom{Shigeru Yamashita}}
    \IEEEauthorblockA{\phantom{Ritsumeikan University} \\
      \phantom{ger@cs.ritsumei.ac.jp}}
    \and
    \IEEEauthorblockN{\phantom{Zanhe Qi}}
    \IEEEauthorblockA{\phantom{Ritsumeikan University} \\ \phantom{goose@ngc.is.ritsumei.ac.jp}}
}
%% \author{
%%     \IEEEauthorblockN{David Clarino}
%%     \IEEEauthorblockA{Ritsumeikan University \\
%%       dclarino@fc.ritsumei.ac.jp}
%%     \and 
%%     \IEEEauthorblockN{Chitranshu Arya}
%%     \IEEEauthorblockA{Netaji Subhas \\
%%       University of Technology \\
%%       chitranshu.arya.ug22@nsut.ac.in}
%%     \and
%%     \IEEEauthorblockN{Shigeru Yamashita}
%%     \IEEEauthorblockA{Ritsumeikan University \\
%%       ger@cs.ritsumei.ac.jp}
%%     \and
%%     \IEEEauthorblockN{Zanhe Qi}
%%     \IEEEauthorblockA{Ritsumeikan University \\
%%       goose@ngc.is.ritsumei.ac.jp}
%% }
\begin{figure*}[t]
  \centering
  \scalebox{1.0} {
    \input{fig-iwz-last}
  }
  \caption{}
  \label{fig-iwz-last}
\end{figure*}
\begin{figure*}[t]
  \centering
  \scalebox{1.0} {
    \input{fig-prop-last}
  }
  \caption{}
  \label{fig-prop-last}
\end{figure*}

\begin{figure}[t]
  \centering
  \scalebox{1.0} {
    \input{xag_and}
  }
  \caption{AND node}
  \label{fig-and}
\end{figure}

\begin{figure}[t]
  \centering
  \scalebox{1.0} {
    \begin{tikzpicture}[
    gate/.style={circle, draw, inner sep=2pt},
    every node/.style={font=\small}
]

% Inputs
\node (x1) at (0,0) {$A$};
\node (x2) at (1.2,0) {$B$};
\node (x3) at (0.6,1.6) {$F$};

% XOR(x1,x2)
\node[gate] (xor12) at (0.6,0.8) {$+$};
\draw (x1) -- (xor12);
\draw (x2) -- (xor12);
\draw[->] (xor12) -- (x3);

\end{tikzpicture}

  }
  \caption{XOR node}
  \label{fig-xor}
\end{figure}

\begin{figure}[t]
  \centering
  \scalebox{1.0} {
    \begin{quantikz}[row sep=0.3cm, column sep=0.55cm]
\qw            & \qw                  & \qw                & \qw   \\
\qw            & \gate[3]{B}          & \gate[3]{A}        & \qw  \\
\qw            & \ghost{B}            & \ghost{A}          & \qw \\
\qw            & \ghost{B}\ctrl{1}    & \ghost{A}\ctrl{1}  & \qw    \\
\qw            & \targ{}              & \targ{}            & \qw  \\
\end{quantikz}
%% \begin{quantikz}[row sep=0.3cm, column sep=0.55cm]
%% \qw            & \qw                            & \gate[3]{A}          & \qw                            & \gate[3]{A}       & \qw   \\
%% \qw            & \gate[3]{B\cdot a_0}         & \ghost{A}            & \gate[3]{B\cdot a_0}         & \ghost{A}         & \qw  \\
%% \qw            & \ghost{B\cdot a_0}           & \ghost{A}\ctrl{1}    & \ghost{B\cdot a_0}           & \ghost{A}\ctrl{1} & \qw \\
%% \lstick{$a_0$} & \ghost{B\cdot a_0}\ctrl{1}   & \targ{}                & \ghost{B\cdot a_0}\ctrl{1}   & \targ{}             & \qw    \\
%% \qw            & \targ{}                        & \qw                    & \targ{}                        & \qw                 & \qw  \\
%% \end{quantikz}


  }
  \caption{XOR diagram}
  \label{fig-xor-qc}
\end{figure}

\begin{figure}[t]
  \centering
  \scalebox{1.0} {
    \input{fig-iz}
  }
  \caption{}
  \label{fig-iz}
\end{figure}

\begin{figure}[t]
  \centering
  \scalebox{1.0} {
    \input{fig-izd}
  }
  \caption{}
  \label{fig-izd}
\end{figure}

\begin{figure}[t]
  \centering
  \scalebox{1.0} {
    %% \begin{quantikz}[row sep=0.25cm, column sep=0.35cm]
%% \lstick{} & \ctrl{2}  & \lstick{} &\ghost{\qw}         & \ghost{\qw} & \qw & \qw                & \targ{}   & \qw      & \ctrl{1} & \gate{T}           & \ctrl{1} & \qw      & \targ{}    & \qw \\
%% \lstick{} & \ctrl{1}  & \lstick{} &\push{=} & \ghost{\qw} & \qw & \qw                & \qw       & \ctrl{1} & \targ{}  & \gate{T^{\dagger}} & \targ{}  & \ctrl{1} & \qw        & \qw \\
%% \lstick{} & \gate{iZ} & \lstick{} &\push{ } & \ghost{\qw} & \qw & \gate{T^{\dagger}} & \ctrl{-2} & \targ{}  & \qw      & \gate{T}           & \qw      & \targ{}  & \ctrl{-2}  & \qw 
%% \end{quantikz}
\begin{quantikz}[row sep=0.25cm, column sep=0.35cm]
\lstick{$\ket{x_a}$} & \ctrl{1}                & \setwiretype{n} &          && \setwiretype{q} & \qw &          & \qw & \qw                & \targ{}   & \qw      & \ctrl{1} & \gate{T}           & \ctrl{1} & \qw      & \targ{}    &          & \qw \\
\lstick{$\ket{x_b}$} & \ctrl{1}                & \setwiretype{n} & \push{=} && \setwiretype{q} & \qw &          & \qw & \qw                & \qw       & \ctrl{1} & \targ{}  & \gate{T^{\dagger}} & \targ{}  & \ctrl{1} & \qw        &          & \qw \\
\lstick{$\ket{x_c}$} & \targ[style={double}]{} & \setwiretype{n} &          && \setwiretype{q} & \qw & \gate{H} & \qw & \gate{T^{\dagger}} & \ctrl{-2} & \targ{}  & \qw      & \gate{T}           & \qw      & \targ{}  & \ctrl{-2}  & \gate{H} & \qw \\
\end{quantikz}


  }
  \caption{}
  \label{fig-rtof}
\end{figure}


\begin{figure}[t]
  \centering
  \scalebox{1.0} {
    \input{fig-iwz-circ}
  }
  \caption{}
  \label{fig-iwz}
\end{figure}

\begin{figure}[t]
  \centering
  \scalebox{1.0} {
    \input{fig-iwz-circd}
  }
  \caption{}
  \label{fig-iwzd}
\end{figure}


\begin{figure}[t]
  \centering
  \scalebox{1.0} {
    \input{fig-iwz-and}
  }
  \caption{}
  \label{fig-iwz-and}
\end{figure}

\begin{figure}[t]
  \centering
  \scalebox{1.0} {
    \input{fig-prop-and}
  }
  \caption{}
  \label{fig-prop-and}
\end{figure}


\begin{algorithm}
\caption{Existing Method}
\begin{algorithmic}[1]

\State $QC \gets$ empty quantum circuit
\State $xag \gets$ XOR-And-Inverter Graph of the function
\State $QC \gets \textsc{parseXAG}(xag)$

\Function{parseXAG}{xag}
    \State $node \gets$ top node of $xag$
    
    \If{$node$ A-child is primary input}
        \State $A \gets$ CNOT with A as control
    \Else
        \State $A \gets \textsc{parseXAG}(\text{subgraph of A-child})$
    \EndIf

    \If{$node$ B-child is primary input}
        \State $B \gets$ CNOT with B as control
    \Else
        \State $B \gets \textsc{parseXAG}(\text{subgraph of B-child})$
    \EndIf

    \If{$node$ is XOR node}
        \State $subQC \gets$ generate circuit from Fig.\ref{fig-xor-qc}
    \ElsIf{$node$ is AND node}
        \If{$A$ and $B$ are primary inputs}
            \State $subQC \gets$ generate circuit Fig.~\ref{fig-iz}
        \ElsIf{node output is overall output}
            \State $subQC \gets$ generate circuit Fig.~\ref{fig-iwz-last}
        \ElsIf{one of $A$, $B$ is primary input}
            \State $subQC \gets$ generate circuit Fig.~\ref{fig-iwz-and}
        \Else
            \State \textbf{error}
        \EndIf
    \EndIf

    \State \Return $subQC$
\EndFunction

\end{algorithmic}
\end{algorithm}

\begin{algorithm}
  \caption{Proposed Method}
  \label{alg-prop}
  \begin{algorithmic}[1]

    \State $QC \gets$ empty quantum circuit
    \State $xag \gets$ XOR-And-Inverter Graph of the function
    \State $QC \gets \textsc{parseXAG}(xag)$

    \Function{parseXAG}{xag}
    \State $node \gets$ top node of $xag$
    
    \If{$node$ A-child is primary input}
    \State $A \gets$ CNOT with A as control
    \Else
    \State $A \gets \textsc{parseXAG}(\text{subgraph of A-child})$
    \EndIf

    \If{$node$ B-child is primary input}
    \State $B \gets$ CNOT with B as control
    \Else
    \State $B \gets \textsc{parseXAG}(\text{subgraph of B-child})$
    \EndIf

    \If{$node$ is XOR node}
    \State $subQC \gets$ generate circuit from Fig.\ref{fig-xor-qc}
    \ElsIf{$node$ is AND node}

    \If{$A$ and $B$ are primary inputs}
       \State $subQC \gets$ generate circuit Fig.~\ref{fig-rtof}
    \ElsIf{node output is overall output}
       \State $subQC \gets$ generate circuit Fig.~\ref{fig-prop-last}
    \ElsIf{one of $A$, $B$ is primary input}
       \State $subQC \gets$ generate circuit Fig.~\ref{fig-prop-and}
    \Else
       \State \textbf{error}
    \EndIf
    \EndIf

    \State \Return $subQC$
    \EndFunction

  \end{algorithmic}
\end{algorithm}


% DCTODO: Add Evaluation, Results, and Conclusion Sections

%\bibliographystyle{unsrt}
%\bibliography{ref}

\end{document}
